% Shared print typography. Content and artwork live in their original sources.
\usepackage[letterpaper,left=0.65in,right=0.65in,top=0.70in,bottom=0.65in,headheight=22pt,headsep=13pt]{geometry}
\usepackage{fontspec,unicode-math}
% A reading face for prose; a contrasting sans face for navigation and diagrams.
\setmainfont{Noto Serif}[Scale=0.96,Ligatures=TeX]
\setsansfont{DejaVu Sans}[Scale=0.92,Ligatures=TeX]
\setmathfont{DejaVu Math TeX Gyre}
\usepackage{graphicx,xcolor,tikz,contour,ragged2e,enumitem,booktabs,longtable,array,tabularx,multicol,needspace,titlesec,fancyhdr,hyperref}
\contourlength{0.45pt}
\definecolor{Ink}{HTML}{17313C}
\definecolor{Teal}{HTML}{087D89}
\definecolor{Gold}{HTML}{A96D16}
\definecolor{Pale}{HTML}{EDF5F5}
\definecolor{Navy}{HTML}{071522}
\definecolor{Muted}{HTML}{536773}
\color{Ink}
\hypersetup{colorlinks=true,linkcolor=Teal,urlcolor=Teal,pdfauthor={Science Olympiad study guides},bookmarksopen=true}
\setlength{\parindent}{0pt}
\setlength{\parskip}{5pt plus1pt minus1pt}
\linespread{1.04}
\setlength{\emergencystretch}{2em}
\raggedbottom
% Ragged-right prose and cells avoid rivers and stretched short descriptions.
\AtBeginDocument{\RaggedRight\lefthyphenmin=3\righthyphenmin=3}
\hyphenpenalty=700
\newcolumntype{Y}{>{\RaggedRight\arraybackslash}X}
% Small labeled, monochrome-safe semantic cues. Never encode meaning by color alone.
% BEGIN GENERATED LEARNING ICONS: shared/learning-icons/generate.py
\newcommand{\cuecanvas}[1]{\tikz[baseline=-1.15em,x=.06em,y=-.06em]{\path[use as bounding box] (0,0) rectangle (24,24);\begin{scope}[line width=.65pt,line cap=round,line join=round]#1\end{scope}}}
\expandafter\def\csname cue@think\endcsname{\cuecanvas{\draw (8,16) .. controls (8,13) and (5,12) .. (5,8) .. controls (5,4) and (8,2) .. (12,2) .. controls (16,2) and (19,4) .. (19,8) .. controls (19,12) and (16,13) .. (16,16) -- cycle;\draw (9,19) -- (15,19) (10,22) -- (14,22) (10,12) -- (12,14) -- (14,12) (12,14) -- (12,16);}}
\expandafter\def\csname cue@worked\endcsname{\cuecanvas{\draw (3,3) -- (10,3) -- (10,8) -- (3,8) -- cycle (14,10) -- (21,10) -- (21,15) -- (14,15) -- cycle (3,17) -- (10,17) -- (10,22) -- (3,22) -- cycle;\draw (10,5) -- (17,5) -- (17,10) (16,15) -- (16,19) -- (10,19);}}
\expandafter\def\csname cue@try\endcsname{\cuecanvas{\draw (4,15) -- (15,4) -- (20,9) -- (9,20) -- (3,21) -- cycle (13,6) -- (18,11) (4,15) -- (9,20);\draw (15,4) -- (17,2) .. controls (18,1) and (19,1) .. (20,2) -- (22,4) .. controls (23,5) and (23,6) .. (22,7) -- (20,9);}}
\expandafter\def\csname cue@review\endcsname{\cuecanvas{\draw (12,12) .. controls (8,10) and (5,10) .. (2,11) -- (2,21) .. controls (5,20) and (8,20) .. (12,22) .. controls (16,20) and (19,20) .. (22,21) -- (22,11) .. controls (19,10) and (16,10) .. (12,12) -- (12,22);\draw (5,7) .. controls (6,3) and (10,1) .. (14,2) .. controls (17,2) and (19,4) .. (20,7) (16,6) -- (20,7) -- (21,3);}}
\expandafter\def\csname cue@?\endcsname{\csname cue@think\endcsname}
\expandafter\def\csname cue@=\endcsname{\csname cue@worked\endcsname}
\expandafter\def\csname cue@>\endcsname{\csname cue@try\endcsname}
\expandafter\def\csname cue@R\endcsname{\csname cue@review\endcsname}
\newcommand{\cueicon}[1]{\csname cue@#1\endcsname}
% END GENERATED LEARNING ICONS
\newcommand{\learningcue}[2]{{\sffamily\small\bfseries\color{Teal}\cueicon{#1}\hspace{.45em}#2}}
\newcommand{\readerguide}{
\subsection*{How to learn with this guide}
Read a short section and connect its explanation to the diagram. Close or cover the text, explain the idea in your own words, then try the question before opening its explanation.
\par\smallskip
\learningcue{think}{Think it through}\quad Explain a cause, make a prediction, or weigh evidence.
\par
\learningcue{worked}{Worked example}\quad Follow the reasoning, then try to reproduce it without looking.
\par
\learningcue{try}{Try it}\quad Apply an idea or method independently.
\par
\learningcue{review}{Review}\quad Recall and connect ideas from earlier sections.
\par\smallskip
Compare your reasoning with the answer, not just your final words or number. Revisit difficult questions on a later day and mix topics instead of only rereading one chapter. Not every lesson needs every kind of cue.
\par
}
\widowpenalty=10000\clubpenalty=10000
\displaywidowpenalty=10000
\setlist[itemize]{leftmargin=1.3em,itemsep=2pt,topsep=3pt}
\setlist[enumerate]{leftmargin=1.5em,itemsep=3pt,topsep=3pt}
\titleformat{\section}{\sffamily\fontsize{20}{24}\selectfont\bfseries\color{Ink}}{}{0pt}{}[\vspace{3pt}{\color{Teal!55}\titlerule[.6pt]}]
\titleformat{\subsection}{\sffamily\fontsize{14}{17}\selectfont\bfseries\color{Ink}}{}{0pt}{}
\titleformat{\subsubsection}{\sffamily\normalsize\bfseries\color{Ink}}{}{0pt}{}
\titlespacing*{\section}{0pt}{19pt plus2pt minus2pt}{10pt}
\titlespacing*{\subsection}{0pt}{16pt plus2pt minus2pt}{6pt}
\titlespacing*{\subsubsection}{0pt}{10pt}{4pt}
\pagestyle{fancy}\fancyhf{}
\newcommand{\GuideChapterMark}{}
\renewcommand{\sectionmark}[1]{\gdef\GuideChapterMark{#1}\markboth{#1}{}}
\fancyhead[L]{\sffamily\fontsize{8}{10}\selectfont\color{Muted}\addfontfeature{LetterSpace=5}SCIENCE OLYMPIAD}
\fancyhead[R]{\sffamily\fontsize{7.2}{9}\selectfont\color{Muted}\parbox[b]{.66\textwidth}{\raggedleft\nouppercase{\leftmark}}}
\fancyfoot[L]{\sffamily\fontsize{8}{10}\selectfont\color{Muted}Solar System B · 2027\\{\fontsize{6.3}{8}\selectfont\sosocredit}}
\fancyfoot[R]{\sffamily\fontsize{8}{10}\selectfont\color{Muted}\thepage}
\renewcommand{\headrulewidth}{0pt}\renewcommand{\footrulewidth}{0pt}
\newcommand{\lesson}[3]{\Needspace{105pt}\subsection{#2}\hypertarget{lesson:#1}{}\label{lesson:#1}{\sffamily\small\color{Teal}#3}\par\nobreak\smallskip}
\newcommand{\subhead}[1]{\par\Needspace{48pt}{\sffamily\fontsize{10.5}{13}\selectfont\bfseries #1}\par\nobreak}
\newcommand{\captiontext}[1]{\par\nobreak\vspace{3pt}{\sffamily\fontsize{9}{11.5}\selectfont\color{Muted}\RaggedRight #1\par}\smallskip}
\newcommand{\checkquestion}[1]{\par\Needspace{50pt}\smallskip\learningcue{think}{Think it through}\quad #1\par\nobreak\smallskip}
\newcommand{\chapterlabel}[1]{\Needspace{180pt}\section{#1}}
\setcounter{secnumdepth}{0}

\newcommand{\workedhead}[1]{\par\Needspace{55pt}\learningcue{worked}{#1}\par\nobreak}

% Measure unbreakable diagram/text rows rather than guessing their height.
% Argument 1 reserves heading space; argument 2 is the heading; argument 3 the row.
\newsavebox{\GuideRowBox}
\newcommand{\measuredrow}[3]{\par\sbox{\GuideRowBox}{\begin{minipage}{\linewidth}#3\end{minipage}}\Needspace{\dimexpr\ht\GuideRowBox+\dp\GuideRowBox+#1pt\relax}\markboth{\GuideChapterMark}{}#2\noindent\usebox{\GuideRowBox}\par\medskip}

% Shared uncluttered cover hierarchy; image-led Solar uses the same type scale.
\newcommand{\manualcover}[6]{\thispagestyle{empty}\vspace*{17pt}
{\sffamily\fontsize{29}{34}\selectfont\bfseries #1}\par\vspace{7pt}
{\sffamily\fontsize{14}{18}\selectfont #2 · 2027}\par\vspace{9pt}
{\sffamily\color{Teal}\large #3}\par\vfill
\begin{center}\includegraphics[width=.85\linewidth,height=330pt,keepaspectratio]{#4}\end{center}
\vfill{\sffamily\large\bfseries #5}\par\smallskip #6\par\vspace*{25pt}\clearpage}

% Full-bleed art option. White type is placed directly on dark artwork, not boxes.
\newcommand{\imageledcover}[6]{\thispagestyle{empty}
\begin{tikzpicture}[remember picture,overlay]
\fill[Navy] (current page.south west) rectangle (current page.north east);
\begin{scope}
\clip (current page.south west) rectangle (current page.north east);
\node[inner sep=0pt] at (current page.center){\includegraphics[width=\paperwidth]{#4}};
\end{scope}
\node[anchor=north west,text width=7.15in,align=left,text=white,inner sep=0pt] at ([xshift=.65in,yshift=-.65in]current page.north west){\sffamily\fontsize{32}{37}\selectfont\bfseries #1\\[7pt]{\fontsize{15}{19}\selectfont #2 · 2027}};
\node[anchor=south west,text width=7.15in,align=left,text=white,inner sep=0pt] at ([xshift=.65in,yshift=.55in]current page.south west){\sffamily\fontsize{17}{21}\selectfont\bfseries #5\\[6pt]{\fontsize{11}{15}\selectfont\mdseries #6}};
\end{tikzpicture}\null\clearpage}
